\begin{figure}[htbp]
  \centering
  \begin{tikzpicture}[x=12mm,y=9mm,alt={A one-dimensional interval runs from zero to seven pixels. The starting point is 6.75, raw 1728, moving right by 1.25, raw 320. Its trial is eight, raw 2048, one pixel beyond the upper wall. Reflecting that excess puts the result at six, raw 1536, with velocity minus 320.}]
    \draw[->] (-.2,0)--(8.5,0);
    \foreach \x in {0,...,8} {\draw (\x,-.1)--(\x,.1);\node[anchor=north,font=\small] at (\x,-.18) {\x};}
    \draw[thick] (7,-.15)--(7,2.5);
    \fill (6.75,0) circle (2pt);
    \draw[->,thick] (6.75,.55)--(8,.55);
    \draw[->,thick] (8,1.25)--(6,1.25);
    \node[anchor=south,font=\small] at (7.5,1.4) {reflect the excess};
    \node[anchor=west,font=\small] at (0,2.2) {$1728+320=2048$};
    \node[anchor=west,font=\small] at (0,1.5) {$3584-2048=1536$};
    \node[anchor=west,font=\small] at (0,.8) {new velocity $=-320$};
  \end{tikzpicture}
  \caption{Reflection preserves the traveled distance after crossing a wall. Coordinates on the axis are pixels; the equations use raw units of one two-hundred-fifty-sixth of a pixel.}
  \label{fig:revision-fixed-reflection}
\end{figure}
